Um die Projektoridentitäten zu beweisen, beginnen wir mit den Definitionen der Projektoren:

\[
P_L = \frac{1 - \gamma^5}{2}, \quad P_R = \frac{1 + \gamma^5}{2}.
\]

Wir nutzen die Eigenschaft von \(\gamma^5\), dass \((\gamma^5)^2 = 1\), um die Identitäten zu beweisen.

Zuerst zeigen wir, dass \(P_L P_R = 0\):

\[
P_L P_R = \left(\frac{1 - \gamma^5}{2}\right)\left(\frac{1 + \gamma^5}{2}\right) = \frac{(1 - \gamma^5)(1 + \gamma^5)}{4} = \frac{1 - (\gamma^5)^2}{4} = \frac{1 - 1}{4} = 0.
\]

Analog zeigt man, dass \(P_R P_L = 0\):

\[
P_R P_L = \left(\frac{1 + \gamma^5}{2}\right)\left(\frac{1 - \gamma^5}{2}\right) = \frac{(1 + \gamma^5)(1 - \gamma^5)}{4} = \frac{1 - (\gamma^5)^2}{4} = 0.
\]

Nun zeigen wir, dass \(P_R P_R = P_R\):

\[
P_R P_R = \left(\frac{1 + \gamma^5}{2}\right)\left(\frac{1 + \gamma^5}{2}\right) = \frac{(1 + \gamma^5)^2}{4} = \frac{1 + 2\gamma^5 + (\gamma^5)^2}{4} = \frac{1 + 2\gamma^5 + 1}{4} = \frac{2 + 2\gamma^5}{4} = \frac{1 + \gamma^5}{2} = P_R.
\]

Ähnlich zeigt man, dass \(P_L P_L = P_L\):

\[
P_L P_L = \left(\frac{1 - \gamma^5}{2}\right)\left(\frac{1 - \gamma^5}{2}\right) = \frac{(1 - \gamma^5)^2}{4} = \frac{1 - 2\gamma^5 + (\gamma^5)^2}{4} = \frac{1 - 2\gamma^5 + 1}{4} = \frac{2 - 2\gamma^5}{4} = \frac{1 - \gamma^5}{2} = P_L.
\]

Nun überprüfen wir, ob die Spinoren \(P_{L, R} u_s(p)\) Eigenzustände des Helizitätsoperators \(\mathfrak{h}\) sind. Der Helizitätsoperator ist gegeben durch:

\[
\mathfrak{h} = \frac{1}{2|\mathbf{p}|}\left(\begin{array}{cc}
\sigma \cdot \mathbf{p} & 0 \\
0 & \sigma \cdot \mathbf{p}
\end{array}\right).
\]

Für masselose Dirac-Spinoren \(u_s(p)\) gilt, dass diese Eigenzustände des Helizitätsoperators mit den Eigenwerten \(\pm \frac{1}{2}\) sind. Dies folgt aus der Tatsache, dass für masselose Teilchen die Helizität eine wohldefinierte Größe ist und die Projektion auf linkshändige oder rechtshändige Komponenten die Helizität nicht ändert.

Für massive Teilchen (\(m \neq 0\)) ist die Helizität nicht mehr eine wohldefinierte Erhaltungsgröße. Um dies zu zeigen, betrachten wir die Transformationseigenschaften unter Lorentztransformationen. Die Helizität ist definiert als

\[
\mathfrak{h} = \frac{\mathbf{S} \cdot \mathbf{p}}{|\mathbf{p}|},
\]

wobei \(\mathbf{S}\) der Spinoperator ist. Bei einer Lorentztransformation ändert sich \(\mathbf{p}\) und damit auch \(\mathfrak{h}\). Da die Masse einen Wechsel zwischen linkshändigen und rechtshändigen Komponenten ermöglicht, sind die Spinoren \(P_{L, R} u_s(p)\) im Allgemeinen keine festen Eigenzustände des Helizitätsoperators mit festen Eigenwerten \(\pm \frac{1}{2}\).

%CHECK: Beweis für den massiven Fall ergänzt, Transformationseigenschaften unter Lorentztransformationen berücksichtigt.